% =============================================================================
% content/experience.tex — Experiência
% =============================================================================

\cvSection{Experiência}

\cvJob{Pesquisador e Desenvolvedor de Inteligência Artificial}
      {(P\&D) Hub de Tecnologia e Inovação (Manaus, Brasil)}
      {2024 – Atualmente}

\cvItem{\textbf{Projetei} e \textbf{implementei} a camada de IA de uma
  plataforma corporativa do zero à produção, desenvolvendo pipelines de RAG e
  backends orientados a eventos (Python, FastAPI, RabbitMQ, MongoDB, Qdrant).}

\cvItem{\textbf{Reduzi} em \textbf{\textasciitilde 50\%} os \textbf{custos
  operacionais} de inferência de LLMs otimizando estratégias de cache e
  filtragem de dados, sem perda mensurável de qualidade.}

\cvItem{Criei APIs REST e bibliotecas (LLMaaS) que expõem modelos de linguagem
  para consumo interno de outras equipes, integrando motores de inferência
  local como \textit{llama.cpp}.}

\cvItem{Desenvolvi integrações resilientes com provedores de IA (OpenAI,
  Anthropic, Together AI), implementando sistemas de fallback, rate limiting e
  tratamento de erros.}

\cvItem{Capacitei agentes autônomos implementando Tool Calling e servidores MCP
  (Model Context Protocol), permitindo que agentes executem ações reais nos
  sistemas da empresa.}

\cvItem{\textbf{Criei uma metodologia própria de avaliação de LLMs}, combinando
  métricas de mercado para mitigar comportamentos indesejados e ranquear
  modelos de forma reprodutível.}

\cvItem{\textbf{Garanti a estabilidade da plataforma} desenvolvendo testes
  unitários (pytest) e ferramentas de teste de carga para simular cenários de
  estresse pré-deploy.}

\cvItem{\textbf{Acelerei o ciclo de entregas} adotando \textit{Spec-Driven
  Development} (SDD) e ferramentas de desenvolvimento assistido por IA (Claude
  Code, Codex), utilizando grafos de conhecimento (\textit{graphify}) para
  manter a precisão do contexto em bases de código extensas.}

\cvItem{\textbf{Atuei como ponte entre tech e business}, traduzindo decisões
  arquiteturais e resultados complexos para stakeholders, além de manter alto
  padrão de código via GitFlow, Code Reviews e registros de decisões técnicas.}

\cvJob{Backend \& iOS Developer}
      {(P\&D) Instituto de Pesquisa Eldorado}
      {2022 – 2024}

\cvItem{Projetei e implementei \textbf{APIs REST} e arquiteturas de
  microsserviços para aplicativos com alto volume de usuários.}

\cvItem{Desenvolvi pipelines de processamento de dados em larga escala
  alimentando sistemas de recomendação baseados em IA.}

\cvItem{Modelei e mantive esquemas de banco relacional e não-relacional
  (\textbf{PostgreSQL}, MongoDB) dando suporte a essas APIs.}

\cvItem{Atuei ativamente em \textbf{code reviews} e documentação técnica,
  elevando o padrão de qualidade e reprodutibilidade do time.}
